% ----------------------------------------------------------------------
%  3.6.6–3.6.9. Магнитное упорядочение, домены, ферромагнетики,
%  антиферро- и ферримагнетики
% ----------------------------------------------------------------------
\subsection{Основы теории магнитного упорядочения}\label{ssec:exchange}

Опыты Эйнштейна--де Хааза и Барнетта (связь намагничивания и вращения
тела) показали, что ферромагнетизм создаётся упорядоченными
\emph{спиновыми} моментами атомов с недостроенными внутренними
оболочками.

\paragraph{Откуда берутся дробные моменты: зоны $s$ и $d$}
У переходного металла уровни $(n-1)d$ геометрически внутренние, но по
энергии лежат выше уровня $ns$. При сближении атомов внешний $ns$-уровень
расщепляется в широкую зону ($N$ подуровней), а $d$-уровни~— в узкую
зону ($5N$ подуровней), лежащую \emph{внутри} $s$-зоны
(рис.~\ref{fig:sd-bands}). Электроны перераспределяются между зонами.

\begin{primer}[Магнитный момент никеля]
Изолированный атом Ni: $3d^84s^2$. Если бы в кристалле $4s$-зона была
занята целиком, а $3d$-зона~— на 8 из 10 мест, остались бы свободные
$3d$-уровни \emph{ниже} занятых $4s$-уровней~— это невыгодно. Часть
$4s$-электронов переходит в $3d$-зону: в металлическом никеле на атом
приходится 9,4 $3d$-электрона и 0,6 $4s$-электрона. До заполнения
$d$-зоны не хватает 0,6 электрона, и некомпенсированный спиновый момент
равен $0{,}6M_\text{Б}$ на атом ($4s$-электроны коллективизированы и в
момент не вкладываются). У железа и кобальта~— $2{,}22M_\text{Б}$ и
$1{,}72M_\text{Б}$. Отсюда и \emph{дробные} моменты.
\end{primer}

\begin{figure}[ht]
\centering
\begin{tikzpicture}
\begin{axis}[width=0.6\textwidth, height=5cm, xmin=0, xmax=1.0, ymin=0, ymax=1.1,
  xlabel={$g(E)$}, ylabel={$E$}, axis lines=left, xtick=\empty, ytick=\empty]
  \addplot[fill=fmBlue!25, draw=none, domain=0:2, samples=241]
    ({x<1 ? 0.14*sqrt(0.86*x)+0.85*sqrt(max(0,1-((0.86*x-0.71)/0.17)^2)) : 0}, {x<1 ? 0.86*x : 0.86*(2-x)});
  \addplot[fmGray, thick, domain=0:1.1, samples=60] ({0.14*sqrt(x)}, {x});
  \addplot[fmRed, thick, domain=0.54:0.88, samples=80] ({0.85*sqrt(max(0,1-((x-0.71)/0.17)^2))+0.14*sqrt(x)}, {x});
  \draw[dashed] (axis cs:0,0.86) -- (axis cs:1,0.86) node[below left, font=\scriptsize]{$\EF$};
  \node[font=\scriptsize, fmRed] at (axis cs:0.55,0.62) {$3d$};
  \node[font=\scriptsize, fmGray, anchor=west] at (axis cs:0.08,0.3) {$4s$};
\end{axis}
\end{tikzpicture}
\caption{Плотность состояний переходного металла (схема для никеля):
узкая высокая $3d$-зона на фоне широкой $4s$-зоны. Уровень Ферми
проходит у самого верха $d$-зоны, оставляя в ней 0,6 незанятого места на
атом. Высокая плотность состояний на уровне Ферми благоприятна для
ферромагнетизма}
\label{fig:sd-bands}
\end{figure}

\paragraph{Обменное взаимодействие}
Перевод спина электрона из антипараллельного положения в параллельное
требует (по принципу Паули) перевода его на более высокий уровень~—
кинетическая энергия растёт, но одновременно меняется и
электростатическая энергия взаимодействия электронов. Суммарное
изменение~— \term{обменная энергия}:
\begin{keyeq}\label{eq:exchange}
E_\text{обм}=-2\sum_{i,j}A_{ij}\,(\vb s_i\vb s_j),
\end{keyeq}
где $A_{ij}$~— \term{обменный интеграл} соседних атомов.
\begin{itemize}
\item $A>0$: выгодны \emph{параллельные} спины~— ферромагнетизм;
\item $A<0$: выгодны \emph{антипараллельные}~— антиферромагнетизм.
\end{itemize}
Параллельные спины вынуждают электроны занимать более высокие уровни
зоны, и ферромагнетизм возникает, только если рост кинетической энергии
меньше выигрыша в обменной энергии. Рост кинетической энергии тем
меньше, чем выше плотность состояний на уровне Ферми,~— поэтому в
выигрышном положении переходные металлы и лантаноиды с узкими $d$- и
$f$-зонами (рис.~\ref{fig:sd-bands}).

Обмен быстро убывает с расстоянием, и в кристалле важен обмен между
ближайшими соседями. Гейзенберг построил теорию, в которой моменты
локализованы у узлов. Часто «магнитные» ионы разделены немагнитными, и
обмен идёт через электроны немагнитного иона (\emph{косвенный обмен});
к ферромагнетизму может приводить и обмен через электроны
проводимости.

\paragraph{Кривая Бете--Слэтера}
Знак и величина обменного интеграла зависят от отношения межатомного
расстояния $a$ к диаметру $2r$ недостроенной $3d$-оболочки
(рис.~\ref{fig:bethe}). Ферромагнитны лишь $\alpha$-Fe, Co и Ni ($a/2r>1{,}5$);
у $\gamma$-Fe, Mn и других $3d$-металлов $a/2r<1{,}5$, и $A<0$. Если
немного «раздвинуть» атомы марганца, можно ожидать ферромагнетизма. Опыт
это подтверждает: ферромагнитны сплавы Гейслера (Cu$_2$MnAl) и соединения
MnSb, MnBi, где атомы марганца удалены друг от друга больше, чем в чистом
марганце; ферромагнетизм появляется и при внедрении в марганец азота,
увеличивающего параметр решётки.

\begin{figure}[ht]
\centering
\begin{tikzpicture}
\begin{axis}[width=0.62\textwidth, height=5cm, xmin=1.38, xmax=2.4, ymin=-1.3, ymax=1.3,
  xlabel={$a/2r$}, ylabel={$A$}, axis lines=middle, ytick=\empty, xtick={1.5},
  every axis x label/.style={at={(axis description cs:1,0.38)}, anchor=north east},
  every axis y label/.style={at={(axis description cs:0,1)}, anchor=south}]
  \addplot[fmBlue, thick, smooth, domain=1.38:2.4, samples=60] {7.2*(x-1.5)*exp(-(x-1.5)/0.35)};
  \fill[fmGray] (axis cs:1.47,-0.235) circle (1.6pt); \node[font=\scriptsize, below right] at (axis cs:1.47,-0.235) {Mn};
  \fill[fmGray] (axis cs:1.42,-0.724) circle (1.6pt); \node[font=\scriptsize, right] at (axis cs:1.42,-0.724) {Cr};
  \fill[fmRed] (axis cs:1.63,0.646) circle (1.6pt); \node[font=\scriptsize, above left] at (axis cs:1.63,0.646) {Fe};
  \fill[fmRed] (axis cs:1.82,0.924) circle (1.6pt); \node[font=\scriptsize, above] at (axis cs:1.82,0.924) {Co};
  \fill[fmRed] (axis cs:1.97,0.883) circle (1.6pt); \node[font=\scriptsize, above right] at (axis cs:1.97,0.883) {Ni};
  \node[font=\scriptsize, fmRed] at (axis cs:2.15,0.25) {ферро};
  \node[font=\scriptsize, fmBlue] at (axis cs:1.45,-1.0) {антиферро};
\end{axis}
\end{tikzpicture}
\caption{Зависимость обменного интеграла от отношения межатомного
расстояния к диаметру $3d$-оболочки (кривая Бете--Слэтера, схема):
ферромагнетизм возможен только при $A>0$}
\label{fig:bethe}
\end{figure}

\begin{glavnoe}
Условия ферромагнетизма: в атоме есть недостроенная \emph{внутренняя}
оболочка ($3d$ или $4f$), а обменный интеграл положителен. Помимо
коллинеарных (параллельных) структур в редкоземельных металлах и их
сплавах встречаются геликоидальные, циклоидальные и синусоидальные
магнитные структуры.
\end{glavnoe}

\paragraph{Точка Кюри}
При нагреве магнитный порядок постепенно разрушается и при точке Кюри
$\theta_K$ (ферро- и ферримагнетики) или Нееля $\theta_N$
(антиферромагнетики) исчезает. Оценим обменный интеграл, приравняв при
$\theta_K$ тепловую энергию обменной (учитывая $z$ ближайших соседей):
$NzA\approx2Nk\theta_K$, откуда
\begin{equation}\label{eq:A-curie}
A\approx\frac{2k\theta_K}z .
\end{equation}
Точки Кюри: Fe~— 1043~К, Co~— 1388~К, Ni~— 627~К, Gd~— 293~К, Dy~— 85~К,
Cu$_2$MnAl~— 630~К, Fe$_3$O$_4$~— 858~К, MnFe$_2$O$_4$~— 560~К.

Как спадает спонтанная намагниченность с температурой, показывает
простейшая теория молекулярного поля Вейсса (рис.~\ref{fig:Ms-T}):
медленно вдали от $\theta_K$ и круто вблизи неё.

\begin{figure}[ht]
\centering
\begin{tikzpicture}
\begin{axis}[width=0.55\textwidth, height=4.8cm, xmin=0, xmax=1.15, ymin=0, ymax=1.08,
  xlabel={$T/\theta_K$}, ylabel={$I_s(T)/I_s(0)$}, axis lines=left, xtick={0,0.5,1}, ytick={0,0.5,1}]
  \addplot[fmRed, thick, domain=0.0005:0.9999, samples=120] ({x/(0.5*ln((1+x)/(1-x)))}, {x});
  \addplot[fmRed, thick] coordinates {(1,0) (1.15,0)};
\end{axis}
\end{tikzpicture}
\caption{Температурная зависимость спонтанной намагниченности в теории
молекулярного поля (расчёт для спина $1/2$: $m=\operatorname{th}(m\,\theta_K/T)$)}
\label{fig:Ms-T}
\end{figure}

\subsection{Доменная структура ферромагнетиков}\label{ssec:domains}

Раз спины в ферромагнетике упорядочены, образец должен быть намагничен
почти до насыщения. На опыте же намагниченность образца часто равна
нулю и растёт до насыщения лишь в поле. Вейсс предположил, что
ферромагнетик разбит на \term{домены}~— области спонтанной
намагниченности с разным её направлением. Их существование доказали
Н.\,С.~Акулов и Ф.~Биттер, а теорию построили Л.\,Д.~Ландау и
Е.\,М.~Лифшиц. Намагниченность образца~— векторная сумма намагниченностей
доменов.

\paragraph{Магнитная анизотропия}
Без поля и напряжений намагниченность домена направлена вдоль
\term{оси лёгкого намагничивания}: в железе (ОЦК)~— $\langle100\rangle$,
в никеле (ГЦК)~— $\langle111\rangle$, в кобальте (ГПУ)~— [0001]. Вдоль лёгких
осей насыщение достигается в слабых полях, вдоль трудных~— лишь в
сильных (рис.~\ref{fig:anisotropy}). Площадь между кривыми лёгкого и
трудного намагничивания~— работа намагничивания вдоль трудной оси,
мера \term{энергии магнитной анизотропии} $K$: для железа при 20\,°C
$K=4{,}2\cdot10^4$~Дж/м$^3$, для кобальта~— $41\cdot10^4$~Дж/м$^3$. В железе
шесть лёгких направлений, в кобальте~— два.

\begin{figure}[ht]
\centering
\begin{tikzpicture}
\begin{axis}[width=0.55\textwidth, height=4.8cm, xmin=0, xmax=1, ymin=0, ymax=1.1,
  xlabel={$H$}, ylabel={$I/I_\text{нас}$}, axis lines=left, xtick=\empty, ytick={1}]
  \addplot[fmRed, thick, domain=0:1, samples=60] {min(1,8*x)};
  \addplot[fmGreen!60!black, thick, domain=0:1, samples=120] {x<0.0884 ? 8*x : 0.707+0.293*min(1,(x-0.0884)/0.35)};
  \addplot[fmBlue, thick, domain=0:1, samples=120] {x<0.0721 ? 8*x : 0.577+0.423*min(1,(x-0.0721)/0.6)};
  \node[font=\scriptsize, fmRed, anchor=west] at (axis cs:0.14,1.04) {[100] лёгкое};
  \node[font=\scriptsize, fmGreen!50!black, anchor=west] at (axis cs:0.5,0.84) {[110]};
  \node[font=\scriptsize, fmBlue, anchor=west] at (axis cs:0.5,0.7) {[111] трудное};
\end{axis}
\end{tikzpicture}
\caption{Кривые намагничивания монокристалла железа вдоль разных
направлений (схема): вдоль лёгкой оси [100] насыщение достигается в
слабом поле}
\label{fig:anisotropy}
\end{figure}

\paragraph{Почему образуются домены}
Рассмотрим одноосный кристалл (кобальт) (рис.~\ref{fig:domain-struct}).
\begin{shagi}
\item Однодоменный кристалл (\textit{а}) создаёт сильное внешнее поле
(свободные полюсы на торцах)~— велика \emph{магнитостатическая} энергия.
\item Разбиение на два антипараллельных домена (\textit{б}) её уменьшает.
\item Замыкающие треугольные домены (\textit{в}) замыкают поток внутри
кристалла, и внешнее поле почти исчезает. Но они намагничены вдоль
трудного направления (энергия анизотропии) и стремятся изменить размеры
(\emph{магнитострикция}), создавая напряжения.
\item В реальной структуре (\textit{г}) объём замыкающих доменов
уменьшается за счёт дробления на тонкие домены, но растёт суммарная
площадь доменных границ, обладающих своей энергией. Равновесие
определяет размер доменов.
\end{shagi}

\begin{figure}[ht]
\centering
\begin{tikzpicture}[scale=0.55, dm/.style={-{Stealth[length=1.6mm]}, line width=0.8pt}]
  % а
  \draw[atomLine] (0,0) rectangle (2.4,3); \draw[dm, fmRed] (1.2,0.4) -- (1.2,2.6);
  \foreach \x in {0.3,0.9,1.5,2.1}{\node[font=\tiny] at (\x,3.25) {N}; \node[font=\tiny] at (\x,-0.25) {S};}
  \node[font=\small] at (1.2,-0.9) {\textit{а}};
  % б
  \begin{scope}[shift={(3.6,0)}]
    \draw[atomLine] (0,0) rectangle (2.4,3); \draw[atomLine, dashed] (1.2,0) -- (1.2,3);
    \draw[dm, fmRed] (0.6,0.4) -- (0.6,2.6); \draw[dm, fmBlue] (1.8,2.6) -- (1.8,0.4);
    \node[font=\small] at (1.2,-0.9) {\textit{б}};
  \end{scope}
  % в
  \begin{scope}[shift={(7.2,0)}]
    \draw[atomLine] (0,0) rectangle (2.4,3);
    \draw[atomLine, dashed] (1.2,0.6) -- (1.2,2.4) (0,0) -- (1.2,0.6) -- (2.4,0) (0,3) -- (1.2,2.4) -- (2.4,3);
    \draw[dm, fmRed] (0.6,0.9) -- (0.6,2.1); \draw[dm, fmBlue] (1.8,2.1) -- (1.8,0.9);
    \draw[dm, fmGreen!60!black] (0.8,2.8) -- (1.6,2.8); \draw[dm, fmGreen!60!black] (1.6,0.2) -- (0.8,0.2);
    \node[font=\small] at (1.2,-0.9) {\textit{в}};
  \end{scope}
  % г
  \begin{scope}[shift={(10.8,0)}]
    \draw[atomLine] (0,0) rectangle (2.4,3);
    \foreach \x in {0.6,1.2,1.8}{\draw[atomLine, dashed] (\x,0.3) -- (\x,2.7);}
    \foreach \x in {0,0.6,1.2,1.8}{\draw[atomLine, dashed] (\x,0) -- (\x+0.3,0.3) -- (\x+0.6,0) (\x,3) -- (\x+0.3,2.7) -- (\x+0.6,3);}
    \draw[dm, fmRed] (0.3,0.5) -- (0.3,2.5); \draw[dm, fmBlue] (0.9,2.5) -- (0.9,0.5);
    \draw[dm, fmRed] (1.5,0.5) -- (1.5,2.5); \draw[dm, fmBlue] (2.1,2.5) -- (2.1,0.5);
    \node[font=\small] at (1.2,-0.9) {\textit{г}};
  \end{scope}
\end{tikzpicture}
\caption{Доменные структуры одноосного кристалла: \textit{а}~—
однодоменная (сильное поле рассеяния); \textit{б}~— два домена;
\textit{в}~— с замыкающими доменами (поток замкнут внутри);
\textit{г}~— реальная структура с мелкими доменами и малыми
замыкающими доменами}
\label{fig:domain-struct}
\end{figure}

\paragraph{Доменная стенка}
Блох показал, что спины не могут резко повернуться на границе доменов:
это стоило бы огромной обменной энергии. Поворот происходит постепенно,
на протяжении многих атомных слоёв~— \term{стенка Блоха}
(рис.~\ref{fig:bloch-wall}). Но чем толще стенка, тем больше в ней
спинов, отклонённых от лёгкой оси, и тем больше энергия анизотропии.
Баланс обменной энергии и энергии анизотропии задаёт толщину стенки~— в
железе около 100~нм. Поскольку стенки обладают энергией, дробление на
домены идёт до определённого предела: домены тем толще, чем больше
обменная энергия и чем более вытянуто тело, и тем тоньше, чем больше
энергия анизотропии.

\begin{masshtab}[ферромагнетизм]{-9.65/-9.45/{обмен между соседями}, -7.2/-6.8/{стенка Блоха}, -6/-3/{домены}}
Обменное взаимодействие действует только между ближайшими соседями
(0,25~нм), но порядок, который оно создаёт, охватывает целый домен.
Стенка Блоха в железе~— около 100~нм (сотни межатомных расстояний),
домены~— от микрометров до долей миллиметра. Частицы мельче примерно
0,1~мкм на домены не делятся и остаются однодоменными.
\end{masshtab}

\begin{figure}[ht]
\centering
\begin{tikzpicture}[x=0.62cm, y=1cm]
  \foreach \i in {0,...,16}{
    \pgfmathsetmacro{\ang}{180*min(1,max(0,(\i-3)/10))}
    \pgfmathsetmacro{\dx}{0.42*sin(\ang)}
    \pgfmathsetmacro{\dy}{0.42*cos(\ang)}
    \pgfmathsetmacro{\shade}{100*min(1,max(0,(\i-3)/10))}
    \draw[-{Stealth[length=1.5mm]}, line width=1pt, fmRed!\shade!fmBlue] (\i-\dx*0.3,-\dy) -- (\i+\dx*0.3,\dy);
    \draw[fmGray!50, line width=0.3pt] (\i,0) ellipse (0.08 and 0.03);}
  \draw[dimline] (3,-0.8) -- node[below, font=\small]{толщина стенки} (13,-0.8);
  \node[font=\small] at (1,0.85) {домен 1};
  \node[font=\small] at (15,0.85) {домен 2};
\end{tikzpicture}
\caption{Стенка Блоха: спин поворачивается на $180^\circ$ постепенно, на
протяжении многих межатомных расстояний}
\label{fig:bloch-wall}
\end{figure}

Очень мелкие частицы на домены не делятся~— выгоднее остаться
\term{однодоменными}: стенка стоила бы слишком много. Предельный размер
однодоменных частиц~— обычно несколько сотен межатомных расстояний.

\paragraph{Намагничивание}
В поле намагниченность растёт двумя путями: \emph{смещением стенок}
(растут домены, намагниченные близко к направлению поля) и \emph{поворотом}
намагниченности доменов к полю. Кривая намагничивания проходит четыре
стадии (рис.~\ref{fig:magn-curve}):
\begin{enumerate}
\item обратимое смещение стенок в слабых полях;
\item необратимое скачкообразное смещение стенок и необратимый поворот
намагниченности доменов к ближайшей к полю лёгкой оси;
\item вращение намагниченности от лёгких осей к направлению поля;
\item насыщение (дальше образец ведёт себя как парамагнетик).
\end{enumerate}
Если поле направлено вдоль лёгкой оси монокристалла, остаются только
стадии 2 и 4, и проницаемость $\mu=1+\d I/\d H$ достигает $10^5$. У
изотропных поликристаллов видны все четыре стадии. Начальная
проницаемость $\mu_\text{нач}$ (стадия~1) важна для слаботочной техники;
с ростом поля проницаемость растёт до $\mu_\text{max}$ (примерно при поле,
равном коэрцитивной силе), а потом спадает к единице.

Дефекты кристалла тормозят смещение стенок; после снятия поля они не дают
стенкам вернуться. Необратимые смещения идут скачками~— \term{эффект
Баркгаузена}: при плавном росте поля намагниченность растёт ступеньками,
и крутой участок кривой состоит из множества мелких скачков.

\begin{figure}[ht]
\centering
\begin{subfigure}[b]{0.48\textwidth}\centering
\begin{tikzpicture}
\begin{axis}[width=\textwidth, height=4.8cm, xmin=0, xmax=1, ymin=0, ymax=1.1,
  xlabel={$H$}, ylabel={$I$}, axis lines=left, xtick=\empty, ytick=\empty]
  \addplot[fmBlue, thick, smooth] coordinates {(0,0) (0.04,0.04) (0.08,0.1) (0.13,0.4) (0.2,0.72) (0.3,0.85) (0.45,0.93) (0.7,0.98) (1,1)};
  \draw[fmGray, dotted] (axis cs:0.075,0) -- (axis cs:0.075,0.1);
  \draw[fmGray, dotted] (axis cs:0.21,0) -- (axis cs:0.21,0.74);
  \draw[fmGray, dotted] (axis cs:0.55,0) -- (axis cs:0.55,0.96);
  \node[font=\scriptsize] at (axis cs:0.035,0.2) {1};
  \node[font=\scriptsize] at (axis cs:0.14,0.82) {2};
  \node[font=\scriptsize] at (axis cs:0.38,0.75) {3};
  \node[font=\scriptsize] at (axis cs:0.78,0.88) {4};
\end{axis}
\end{tikzpicture}
\caption{}
\end{subfigure}\hfill
\begin{subfigure}[b]{0.48\textwidth}\centering
\begin{tikzpicture}
\begin{axis}[width=\textwidth, height=4.8cm, xmin=-1.1, xmax=1.1, ymin=-1.15, ymax=1.15,
  axis lines=middle, xtick={-0.3}, xticklabels={$-H_\text{к}$}, ytick={0.834}, yticklabels={$I_\text{ост}$},
  xlabel={$H$}, ylabel={$I$}, xticklabel style={font=\scriptsize}, yticklabel style={font=\scriptsize}]
  \addplot[fmRed, thick, domain=-1.1:1.1, samples=80] {tanh(4*(x+0.3))};
  \addplot[fmRed, thick, domain=-1.1:1.1, samples=80] {tanh(4*(x-0.3))};
  \addplot[fmBlue, thick, dashed, domain=0:1.1, samples=40] {tanh(3.2*x)^2*sign(x)};
\end{axis}
\end{tikzpicture}
\caption{}
\end{subfigure}
\caption{\textit{а}~— кривая намагничивания ферромагнетика и её четыре
стадии (схема); \textit{б}~— петля гистерезиса: $I_\text{ост}$~—
остаточная намагниченность, $H_\text{к}$~— коэрцитивная сила (пунктир~—
начальная кривая намагничивания)}
\label{fig:magn-curve}
\end{figure}

\subsection{Магнитные свойства ферромагнетиков}\label{ssec:ferro}

Ферромагнитны девять элементов: три $3d$-металла (Fe, Co, Ni) и шесть
тяжёлых редкоземельных $4f$-металлов (Gd, Tb, Dy, Ho, Er, Tm), а также
огромное число сплавов и соединений с участием элементов с недостроенными
$d$- и $f$-оболочками.

\paragraph{Характеристики}
Ферромагнитный материал описывают несколькими величинами; одни
определяются только составом, другие~— структурой.
\begin{itemize}
\item \term{Намагниченность насыщения} $I_\text{нас}$ определяется лишь
числом некомпенсированных спинов: для железа при 0~К
$I_\text{нас}=2{,}22M_\text{Б}N$, для никеля~— $0{,}6M_\text{Б}N$ ($N$~— число
атомов в единице объёма). От структуры она почти не зависит.
\item \term{Гистерезис} (рис.~\ref{fig:magn-curve},\,\textit{б}): при
уменьшении поля намагниченность не возвращается по исходной кривой. При
$H=0$ остаётся \term{остаточная намагниченность} $I_\text{ост}$ (индукция
$B_\text{ост}$); чтобы её снять, нужно обратное поле~— \term{коэрцитивная
сила} $H_\text{к}$. Площадь петли пропорциональна работе
перемагничивания за цикл. При малой амплитуде поля, когда работает
только обратимое смещение стенок, петля вырождается в прямую.
\item $I_\text{ост}$~— структурно чувствительна: зависит от текстуры и
направления поля относительно лёгких осей. Для изотропного материала,
перемагничиваемого поворотом доменов, $I_\text{ост}\approx0{,}5I_\text{нас}$.
\item $H_\text{к}$~— тоже структурно чувствительна и определяется главным
образом дефектами: торможением стенок в поле внутренних напряжений и
закреплением их на немагнитных включениях. Сильнее всего закрепляют
стенку включения с размером порядка её толщины.
\end{itemize}

\paragraph{Деформация и нагрев}
Холодная деформация создаёт внутренние напряжения, затрудняющие
намагничивание: проницаемость падает, коэрцитивная сила растёт тем
больше, чем сильнее обжатие. Рекристаллизационный отжиг возвращает
свойства к исходным. При нагреве намагниченность сначала медленно, а
затем резко падает и в точке Кюри практически исчезает (остаётся слабая
парамагнитная). Проницаемость в сильных полях с нагревом уменьшается, а
в слабых сначала \emph{растёт} (ослабевают силы, мешающие смещению
стенок), а у точки Кюри резко падает (исчезает спонтанная намагниченность);
выше $\theta_K$ проницаемость близка к единице\ispr{В оригинале (с.~553)
сказано, что проницаемость в точке Кюри «падает до нуля»; падает до
нуля восприимчивость (точнее, до малых парамагнитных значений), а
проницаемость $\mu=1+\chi$~— до значений, близких к единице.}.

\paragraph{Магнитные материалы}
По сочетанию свойств материалы делят на две группы (табл.~\ref{tab:magn-materials}).
\begin{itemize}
\item \term{Магнитомягкие}: намагничиваются до насыщения в слабых полях,
высокая начальная проницаемость, малые потери на перемагничивание;
намагничивание идёт в основном смещением стенок; условно
$H_\text{к}<800$~А/м. Применение~— сердечники трансформаторов,
электромагнитов, магнитные экраны.
\item \term{Магнитотвёрдые}: высокая коэрцитивная сила ($H_\text{к}>4$~кА/м)
и большая запасённая энергия; намагничивание в основном вращением
(в некоторых РЗМ-сплавах~— смещением стенок). Применение~— постоянные
магниты и носители информации.
\end{itemize}

\begin{table}[ht]
\centering
\caption{Магнитомягкие и магнитотвёрдые материалы}
\label{tab:magn-materials}
\small
\begin{tabularx}{\textwidth}{@{}>{\raggedright\arraybackslash}X>{\raggedright\arraybackslash}X@{}}
\toprule
Магнитомягкие ($H_\text{к}<800$~А/м) & Магнитотвёрдые ($H_\text{к}>4$~кА/м) \\
\midrule
технически чистое железо; электротехническая сталь Fe--(0,5--5\,\%)Si;
сендаст Fe--9,5\,\%Si--5,5\,\%Al; пермаллои Fe--(50--80\,\%)Ni; аморфные
сплавы на основе Fe и Co; ферриты MO$\cdot$Fe$_2$O$_3$ (M = Ni, Zn, Mn) &
сплавы Fe--Ni--Al(--Co) (альнико); Nd--Fe--B; SmCo$_5$; гексаферриты
бария и стронция MO$\cdot$6Fe$_2$O$_3$ (M = Ba, Sr); феррит кобальта
CoO$\cdot$Fe$_2$O$_3$\ispr{В оригинале (с.~554) магнитотвёрдые ферриты
записаны общей формулой MO$\cdot$Fe$_2$O$_3$ (M = Ba, Co, Sr); бариевый
и стронциевый ферриты имеют состав MO$\cdot$6Fe$_2$O$_3$ (гексаферриты).} \\
\bottomrule
\end{tabularx}
\end{table}

\paragraph{Магнитострикция}
При намагничивании ферромагнетик меняет размеры и форму~—
\term{магнитострикция}. Пусть малое однодоменное тело выше точки Кюри~—
шар. При охлаждении ниже $\theta_K$ межатомное расстояние меняется так,
чтобы обменный интеграл стал выгоднее: у железа и кобальта (возрастающий
участок кривой Бете--Слэтера) объём растёт, у никеля (спадающий участок)
уменьшается. Кроме того, из-за взаимодействия выстроенных вдоль лёгкой
оси моментов шар становится эллипсоидом. В многодоменном кристалле
соседние домены мешают друг другу менять форму, и возникают упругие
напряжения.

\paragraph{Ферромагнетизм сплавов}
Средний атомный момент ферромагнетиков нецелый: Fe~— 2,22, Co~— 1,72,
Ni~— 0,60~$M_\text{Б}$. В сплавах он меняется, и проще всего это понять
на никеле (модель жёсткой зоны). В $d$-зоне никеля не хватает 0,6
электрона на атом.
\begin{itemize}
\item Медь ($3d^{10}4s^1$) приносит на один электрон больше, чем никель; он
попадает в $3d$-зону и заполняет её. Момент падает с концентрацией меди
и при 60~ат.\,\% Cu обращается в нуль (рис.~\ref{fig:slater-ni}).
\item Цинк (на два электрона больше) уменьшает момент вдвое быстрее,
алюминий (на три)~— втрое.
\item Палладий имеет такое же внешнее строение, как никель, и момента не
меняет.
\end{itemize}
Сходные закономерности найдены для сплавов Fe, Co, Ni между собой и с
другими $3d$-элементами; число некомпенсированных спинов в $d$-зоне
таких сплавов не превышает 2,4 на атом.

\begin{figure}[ht]
\centering
\begin{tikzpicture}
\begin{axis}[width=0.58\textwidth, height=4.8cm, xmin=0, xmax=65, ymin=0, ymax=0.7,
  xlabel={концентрация добавки, ат.\,\%}, ylabel={$\mu/M_\text{Б}$}, axis lines=left,
  yticklabel style={/pgf/number format/use comma}]
  \addplot[fmGray, thick] coordinates {(0,0.6) (60,0.6)};
  \addplot[fmBlue, thick] coordinates {(0,0.6) (60,0)};
  \addplot[fmGreen!60!black, thick] coordinates {(0,0.6) (30,0)};
  \addplot[fmRed, thick] coordinates {(0,0.6) (20,0)};
  \node[font=\scriptsize, fmGray, anchor=south] at (axis cs:45,0.6) {Pd};
  \node[font=\scriptsize, fmBlue, anchor=west] at (axis cs:42,0.2) {Cu};
  \node[font=\scriptsize, fmGreen!50!black, anchor=west] at (axis cs:21,0.2) {Zn};
  \node[font=\scriptsize, fmRed, anchor=east] at (axis cs:13,0.15) {Al};
\end{axis}
\end{tikzpicture}
\caption{Средний магнитный момент атома в твёрдых растворах на основе
никеля (модель жёсткой зоны): каждый лишний валентный электрон добавки
заполняет $d$-зону никеля}
\label{fig:slater-ni}
\end{figure}

Точку Кюри двойных неупорядоченных растворов рассчитывают по теории
Гейзенберга--Блоха--Вонсовского через обменные интегралы пар A--A, B--B,
A--B и концентрации $c_A+c_B=1$:
\begin{equation}\label{eq:curie-alloy}
\theta_K=\frac z{2k}\bigl(c_A^2A_{AA}+c_B^2A_{BB}+c_Ac_BA_{AB}\bigr);
\end{equation}
расчёт хорошо согласуется с опытом (растворы Co, Fe, Pd, Pt в никеле).
Растворение неферромагнитных компонентов почти всегда понижает точку
Кюри.

\subsection{Антиферромагнетики и ферримагнетики}\label{ssec:antiferro}

При $A<0$ спины соседей антипараллельны. Если моменты соседних узлов
равны~— \term{антиферромагнетизм}, если нет~— \term{ферримагнетизм}.
Антиферромагнитны хром (до 311~К), марганец (ниже $\approx100$~К),
лёгкие редкоземельные металлы (Ce, Pr, Nd, Sm, Eu~— при низких
температурах)\ispr{В оригинале (с.~557) эти металлы названы
антиферромагнетиками при комнатной температуре. Это так лишь для
хрома ($\theta_N\approx311$~К); у марганца $\theta_N\approx95$--100~К, у
лёгких РЗМ~— от долей до десятков кельвинов (Eu~— около 90~К).} и
множество соединений переходных металлов. Восприимчивость
антиферромагнетика растёт с температурой до максимума в \term{точке
Нееля} $\theta_N$, где порядок исчезает, а выше уменьшается, как у
парамагнетика.

\paragraph{Магнитная структура}
Магнитную структуру~— расположение моментов в решётке~— изучают по
дифракции нейтронов (у нейтрона есть магнитный момент). Антиферромагнетик
можно представить как две (или более) вставленные друг в друга
подрешётки с параллельными моментами, компенсирующие друг друга. Пример~—
оксид марганца MnO со структурой NaCl (рис.~\ref{fig:MnO}): соседние
плоскости \{111\} ионов Mn$^{2+}$ намагничены противоположно, и период
магнитной структуры вдвое больше периода атомной. Магнитное
взаимодействие плоскостей слегка искажает кубическую решётку, что видно
рентгенографически.

\begin{figure}[ht]
\centering
\begin{tikzpicture}[cubeproj, scale=2.0]
  \CubeEdges
  \foreach \i in {0,1,2}{\foreach \j in {0,1,2}{\foreach \k in {0,1,2}{
    \pgfmathtruncatemacro{\s}{\i+\j+\k}
    \pgfmathtruncatemacro{\ev}{mod(\s,2)}
    \ifnum\ev=0
      \pgfmathtruncatemacro{\pl}{mod(\s/2,2)}
      \coordinate (P) at (\i/2,\j/2,\k/2);
      \ifnum\pl=0
        \draw[-{Stealth[length=1.6mm]}, fmRed, line width=1pt] ($(P)+(0,0,-0.11)$) -- ($(P)+(0,0,0.11)$);
      \else
        \draw[-{Stealth[length=1.6mm]}, fmBlue, line width=1pt] ($(P)+(0,0,0.11)$) -- ($(P)+(0,0,-0.11)$);
      \fi
    \fi}}}
\end{tikzpicture}
\caption{Антиферромагнитная структура MnO (показаны только ионы
Mn$^{2+}$, образующие ГЦК-подрешётку; направление стрелок условно): ионы
каждой плоскости \{111\} намагничены одинаково, соседние плоскости~—
противоположно}
\label{fig:MnO}
\end{figure}

\paragraph{Ферриты}
Самые типичные ферримагнетики~— \term{ферриты} MeO$\cdot$Fe$_2$O$_3$ (Me~—
Mg$^{2+}$, Zn$^{2+}$, Cu$^{2+}$, Ni$^{2+}$, Fe$^{2+}$, Mn$^{2+}$) со
структурой \emph{шпинели}. В элементарной ячейке 8 формульных единиц: 32
иона кислорода в плотной кубической упаковке, 8 ионов Me$^{2+}$ и 16 ионов
Fe$^{3+}$ в её пустотах. В распространённых \emph{обращённых} шпинелях
(например, Fe$_3$O$_4$, NiFe$_2$O$_4$) 8 ионов Fe$^{3+}$ занимают
тетраэдрические пустоты, а 8 ионов Me$^{2+}$ и 8 ионов Fe$^{3+}$~—
октаэдрические\ispr{В оригинале (с.~558) сказано, что ионы Fe$^{3+}$
занимают 8 октаэдрических пустот, а 8 Me$^{2+}$ и 8 Fe$^{3+}$~— 16
тетраэдрических. В шпинели наоборот: заняты 8 тетраэдрических (A) и 16
октаэдрических (B) позиций; в обращённой шпинели в A~— Fe$^{3+}$, в B~—
Me$^{2+}$ и Fe$^{3+}$.}. Моменты подрешёток A и B антипараллельны:
моменты ионов Fe$^{3+}$ в них взаимно компенсируются, и наблюдаемый момент
создают ионы Me$^{2+}$.

\begin{primer}[Магнетит]
В Fe$_3$O$_4$=FeO$\cdot$Fe$_2$O$_3$ моменты двух ионов Fe$^{3+}$ ($5M_\text{Б}$
каждый) компенсируются, и момент формульной единицы создаёт ион Fe$^{2+}$~—
$4M_\text{Б}$. На один атом железа~— $4/3\approx1{,}33M_\text{Б}$, что близко к
опыту.
\end{primer}

Ферриты~— ионные кристаллы с плохой электропроводностью. В сочетании с
хорошими магнитными свойствами это сделало их незаменимыми в технике
высоких и сверхвысоких частот, где металлические ферромагнетики
непригодны из-за потерь на вихревые токи.

Магнитное упорядочение найдено и в аморфных телах. Есть и
кристаллические вещества с хаотически «замороженными» спинами~—
\term{спиновые} (магнитные) \term{стёкла}.

\begin{itogi}
\begin{itemize}
\item $\vb B=\mu_0(\vb H+\vb I)$, $\vb I=\chi\vb H$; диа- и парамагнетики~—
слабомагнитные ($|\chi|\sim10^{-7}$--$10^{-5}$), ферро- и ферримагнетики~—
сильномагнитные ($\chi$ до $10^5$).
\item Магнетон Бора $M_\text{Б}=e\hbar/2m$; магнитны атомы с недостроенными
$d$- и $f$-оболочками (правило Хунда).
\item Диамагнетизм~— ларморова прецессия, $\chi<0$, не зависит от $T$.
Парамагнетизм~— ориентация моментов, $\chi=C/T$ (Кюри) или $C/(T-\theta)$
(Кюри--Вейсс); электроны проводимости~— парамагнетизм Паули, не
зависящий от $T$.
\item Ферромагнетизм: обменный интеграл $A>0$ (кривая Бете--Слэтера);
дробные моменты из-за $s$--$d$-перераспределения электронов; точка Кюри
$\sim A z/2k$.
\item Домены уменьшают магнитостатическую энергию; стенка Блоха~—
компромисс обмена и анизотропии; намагничивание идёт смещением стенок и
вращением; гистерезис, $I_\text{ост}$, $H_\text{к}$~— структурно
чувствительны.
\item Антиферромагнетики~— компенсированные подрешётки (точка Нееля);
ферриты~— некомпенсированные подрешётки шпинели.
\end{itemize}
\end{itogi}

\begin{proverka}
\item Почему диамагнетизм присущ всем веществам, а наблюдается в чистом
виде лишь у некоторых?
\item Почему восприимчивость Паули не зависит от температуры, а
восприимчивость парамагнитной соли подчиняется закону Кюри?
\item Объясните, почему магнитный момент атома никеля в металле равен
$0{,}6M_\text{Б}$.
\item Почему марганец не ферромагнитен, а сплавы Гейслера на его основе
ферромагнитны?
\item Зачем ферромагнетику домены? Почему мелкие частицы однодоменны?
\item Почему наклёп увеличивает коэрцитивную силу?
\item При какой концентрации цинка исчезнет ферромагнетизм никеля в модели
жёсткой зоны?
\end{proverka}
